% !TEX root = ../main.tex
%==============================================================
\chapter{Kết luận và hướng phát triển}
\label{ch:ket-luan}

\section{Kết luận}
Đề tài đã nghiên cứu và xây dựng thành công một bộ kết xuất Path Tracing chạy
song song trên kiến trúc CUDA của NVIDIA, đạt được các mục tiêu đề ra ở
Chương~\ref{ch:gioi-thieu}. Cụ thể, hệ thống đã:
\begin{itemize}[noitemsep]
  \item Hiện thực trọn vẹn thuật toán Path Tracing trên GPU bằng CUDA, với vòng
        dò đường theo kiểu lặp và cơ chế tích lũy mẫu theo thời gian để khử
        nhiễu Monte Carlo.
  \item Xây dựng hệ thống vật liệu Lambertian, kim loại và điện môi cùng mô hình
        camera vật lý hỗ trợ độ sâu trường ảnh và nhòe chuyển động.
  \item Hiển thị thời gian thực nhờ cầu nối CUDA--OpenGL, tránh chi phí sao chép
        dữ liệu ảnh qua lại giữa CPU và GPU.
  \item Nạp cảnh theo định dạng Mitsuba~3 và cung cấp giao diện điều khiển tương
        tác để tinh chỉnh camera cũng như theo dõi hiệu năng.
\end{itemize}

Kết quả thực nghiệm ở Chương~\ref{ch:danh-gia} xác nhận hệ thống tái tạo đầy đủ
các hiệu ứng quang học mục tiêu và duy trì tốc độ tương tác trên phần cứng phổ
thông. Đồng thời, thực nghiệm cũng chỉ ra nút thắt hiệu năng khi mật độ vật thể
tăng cao, do việc tìm giao điểm còn thực hiện bằng duyệt tuyến tính, mở ra các
hướng phát triển tiếp theo.

\section{Hướng phát triển}
Trên nền tảng hiện có, hệ thống có thể được mở rộng theo các hướng sau:
\begin{enumerate}[noitemsep]
  \item \textbf{Bounding Volume Hierarchy (BVH):} xây dựng cấu trúc tăng tốc
        không gian để giảm chi phí tìm giao điểm từ tuyến tính xuống mức
        logarit, khắc phục nút thắt hiệu năng ở cảnh mật độ cao.
  \item \textbf{1D/2D:} bổ sung các loại hình học và nguyên thủy một, hai chiều
        (đoạn thẳng, mặt phẳng, tam giác) bên cạnh khối cầu hiện tại.
  \item \textbf{Texture Mapping:} hỗ trợ ánh xạ vân bề mặt để tăng độ chi tiết
        và chân thực cho vật liệu.
  \item \textbf{Emission Material:} thêm vật liệu phát sáng để mô phỏng nguồn
        sáng vùng trực tiếp trong cảnh.
  \item \textbf{Volume Rendering:} kết xuất môi trường tham gia (khói, sương,
        chất bán trong suốt) bằng mô hình tán xạ thể tích.
  \item \textbf{Monte Carlo:} nâng cao chất lượng ước lượng Monte Carlo bằng lấy
        mẫu theo tầm quan trọng (Importance Sampling) và ước lượng sự kiện kế
        tiếp để giảm nhiễu nhanh hơn.
\end{enumerate}

\section{Mã nguồn}
Toàn bộ mã nguồn của đề tài được công khai tại kho lưu trữ:
\begin{center}
\url{https://github.com/KinhLupNe/Path-tracing-Render-CUDA}
\end{center}
